

\textbf{OpenEvolve-MARL}

Open-Source Adaptive Evolutionary Algorithm Discovery for Multi-Agent
Systems

Project Proposal \textbar{} Compute Budget: 8 × H100 GPUs

\textbf{1. Theme and Motivation}

The central challenge in multi-agent reinforcement learning (MARL) ---
particularly in imperfect-information game solving --- has long been
that state-of-the-art algorithmic variants such as Discounted CFR (DCFR)
and SHOR-PSRO emerge not from systematic search, but from human
intuition applied painstakingly over years. The recent paper AdaEvolve
(Cemri et al., 2026, arXiv:2602.20133) demonstrates that this bottleneck
can be broken: a hierarchical adaptive optimization framework wrapped
around large language model (LLM) mutation operators can discover novel
algorithms --- autonomously, at scale, and with greater compute
efficiency than prior systems like AlphaEvolve. However, AdaEvolve, like
its predecessors, is built entirely on proprietary models (e.g., Gemini
Flash / Pro). This creates a fundamental reproducibility crisis and a
hard dependency on API cost and rate limits that makes democratized
scientific discovery essentially impossible.

OpenEvolve-MARL proposes to replicate and extend
AdaEvolve\textquotesingle s core architecture using exclusively
open-source LLMs --- specifically models from the Qwen 2.5 Coder,
DeepSeek-Coder-V3, and Llama 3.3 families --- running on our 8 × H100
cluster. Beyond simple replication, we will use the evolutionary loop
itself as a training signal to fine-tune the mutation LLMs via
reinforcement learning from execution feedback (RLEF), creating a
self-improving discovery engine. The domain of application is the same
as in the original paper: evolving CFR variants for regret minimization
and meta-solver logic for PSRO, with the explicit goal of exceeding the
performance of both human-designed and proprietary-LLM-discovered
baselines.

The timing is ideal. The open-source ecosystem now has coding LLMs
competitive with GPT-4o on standard benchmarks, quantization tooling
(AWQ, GPTQ) that allows 70B-parameter models to run at high throughput
on 8 × H100, and efficient serving frameworks (vLLM, SGLang) that make
population-scale evolutionary sampling tractable without proprietary
APIs. This project is designed to sit at the intersection of three
high-impact frontiers: automated algorithm discovery, open-source AI
research infrastructure, and MARL theory.

\textbf{2. Method}

\textbf{2.1 Framework Architecture}

We will build directly on the AdaEvolve three-level hierarchy --- Local
Adaptation (exploration intensity per island), Global Adaptation
(bandit-based compute routing across islands), and Meta-Adaptation
(dynamic prompt generation from an accumulated improvement signal) ---
but replace every proprietary model call with local inference. The
architecture will consist of two LLM tiers analogous to
AlphaEvolve\textquotesingle s Flash/Pro split:

\begin{itemize}
\item
  High-throughput Mutation Agent: Qwen2.5-Coder-7B-Instruct or
  DeepSeek-Coder-1.3B, quantized to INT4 via AWQ, served with vLLM for
  maximum token throughput (\textasciitilde2,000 tokens/sec on a single
  H100). This agent handles the bulk of population mutations.
\item
  High-quality Proposal Agent: DeepSeek-Coder-V3 (7B parameter
  distillate) or Qwen2.5-Coder-32B, served at FP8 precision. This agent
  is invoked selectively by the global bandit for high-promise islands,
  generating more structurally novel mutations.
\end{itemize}

The evolutionary search targets the same evolvable code skeletons as in
the original paper: (i) the update\_accumulate\_regret,
get\_updated\_current\_policy, and update\_accumulate\_policy methods
for CFR, and (ii) the TrainMetaStrategySolver and EvalMetaStrategySolver
for PSRO. These skeletons define the search space while keeping the
fitness evaluator (exploitability computation on Kuhn Poker, Leduc
Hold\textquotesingle em, and Blotto) deterministic and fast.

\textbf{2.2 Fine-Tuning Pipeline}

The most technically novel contribution of this project is closing the
loop between evolutionary search and model improvement. After each
generation of the evolutionary search, any mutation that results in a
fitness improvement above a threshold delta becomes a positive training
example: the (prompt, diff) pair is retained for fine-tuning. Mutations
that compile and run but worsen fitness provide hard negatives. We will
apply Direct Preference Optimization (DPO) with a short rolling window
of examples to continuously update the high-throughput mutation agent.
This is feasible on our hardware: with LoRA adapters (rank 16, targeting
q/v projections), a DPO update over 512 examples on a 7B model takes
roughly 8 minutes on a single H100, meaning we can interleave one
fine-tuning pass per every 20--30 evolutionary generations without
meaningfully interrupting the search. The expected effect is that the
mutation agent internalizes domain-specific code patterns --- CFR update
structures, regret discount schemes --- and produces higher-quality
proposals over time, a dynamic not present in any prior AdaEvolve-style
system.

\textbf{2.3 Compute Allocation}

Across our 8 × H100 cluster we propose the following allocation. Six
GPUs will run the vLLM inference server hosting both agent tiers,
partitioned so that four cards serve the high-throughput agent (tensor
parallel 2) and two cards serve the high-quality agent. One GPU will run
the evolutionary orchestrator, fitness evaluators (which are CPU-light
but benefit from GPU for batch Nash computation at scale), and the
program database. The final GPU is reserved for DPO fine-tuning runs. A
complete evolutionary run targeting CFR on Leduc Hold\textquotesingle em
for 1,000 generations is estimated to consume approximately 40--60
GPU-hours, meaning we can run multiple full experiments, ablations, and
domain extensions within a reasonable wall-clock budget.

\textbf{3. Evaluation Plan}

We will evaluate on three dimensions: algorithmic performance, discovery
efficiency, and model improvement signal. For algorithmic performance,
the primary metric is exploitability (measured in milli-big-blinds for
poker games) of evolved CFR variants versus DCFR, PCFR+, and the
AdaEvolve-discovered VAD-CFR baseline. For PSRO, we track NashConv
against SHOR-PSRO on normal-form games. We will also measure wall-clock
convergence time at fixed exploitability thresholds, since MARL
practitioners care about iteration speed as much as asymptotic
performance.

For discovery efficiency, we will report the number of LLM calls
required to first exceed each baseline --- directly comparable to the
numbers reported in AdaEvolve --- as well as the fraction of generated
mutations that compile and improve fitness (the \textquotesingle useful
mutation rate\textquotesingle). This metric is where we expect the DPO
fine-tuning to show the most measurable gain over time, and where our
results will most directly demonstrate the value of the self-improving
loop relative to a frozen open-source model. Finally, we will run an
ablation comparing our two-tier open-source ensemble against a
single-model configuration and against the AdaEvolve setup, to isolate
the contribution of the tiered architecture independent of model
quality.

\textbf{4. Risks and Mitigations}

The primary technical risk is that open-source coding LLMs produce a
substantially lower useful mutation rate than Gemini Pro, leading to
stagnation in the evolutionary search. Our mitigation is twofold: the
DPO fine-tuning loop is specifically designed to address this over time,
and the AdaEvolve hierarchical adaptation mechanism will reduce wasted
compute on stagnating islands regardless of per-mutation quality. A
secondary risk is that the DPO signal is too sparse or noisy early in
training, before the population has produced enough positive examples.
We will address this with a warm-start phase using synthetic preference
data generated by prompting the high-quality agent to rate and rewrite
its own outputs. A third risk is that inference throughput on the
quantized 7B model is insufficient to sustain a competitive evolutionary
rate. We will validate throughput in the first week and fall back to a
smaller model (1.3B) if necessary, accepting a trade-off in mutation
quality. Finally, there is the question of fitness evaluator coverage:
exploitability is only tractable for small games. For larger games we
will use approximate best-response computation with a fixed oracle
budget, which introduces noise. We will report confidence intervals and
validate trends on small games first.

\textbf{5. Potential Impact}

This project\textquotesingle s potential impact extends well beyond the
specific domain of imperfect-information game solving, and we believe it
constitutes one of the highest-leverage uses of 8 × H100 GPUs currently
available to an academic or independent research team.

\textbf{5.1 Democratizing Automated Scientific Discovery}

AlphaEvolve and AdaEvolve have demonstrated that LLM-driven evolutionary
search can discover algorithms competitive with or superior to decades
of human expert work. But every result to date has been produced using
proprietary models behind closed APIs. This means the community cannot
reproduce the results, cannot inspect the models, cannot fine-tune them
on domain-specific data, and cannot run experiments without per-token
costs that are prohibitive for most research groups. OpenEvolve-MARL
would be the first fully reproducible, open-weight implementation of an
adaptive LLM-driven evolutionary discovery system. If successful, it
becomes infrastructure --- a platform other research groups can deploy
on their own hardware to run evolutionary searches in their own domains.
The multiplier effect on the broader community could be enormous: every
domain in which AdaEvolve-style search could discover better algorithms
(numerical optimization, compiler passes, cryptographic primitives, ML
training recipes) becomes accessible to researchers who currently cannot
afford proprietary API costs at the scale required.

\textbf{5.2 Self-Improving Discovery as a Research Paradigm}

The DPO fine-tuning loop is not merely an engineering improvement over
AdaEvolve --- it represents a qualitatively new paradigm. In standard
evolutionary search, the mutation operator is stationary: its quality is
fixed by pretraining. In our system, the mutation operator improves as
it accumulates evidence about what kinds of code changes produce fitness
gains in the target domain. This is, in essence, a form of online
meta-learning for algorithm discovery. If the fine-tuning signal is
strong enough, the mutation LLM will converge toward a specialist that
generates domain-relevant code perturbations far more reliably than a
general-purpose coding model. This has implications beyond MARL: the
same architecture could be applied to any domain with a fast,
differentiable fitness evaluator, and the resulting specialist model
would be a more efficient discovery engine than a frontier proprietary
model for that specific domain. Demonstrating this principle empirically
--- even on a small scale, on 8 × H100 --- would be a significant
contribution to the growing literature on inference-time compute scaling
and self-improving AI systems.

\textbf{5.3 Direct Impact on Multi-Agent Systems Research}

Within MARL specifically, the algorithms discovered by AdaEvolve
(VAD-CFR, SHOR-PSRO) have already outperformed state-of-the-art
baselines on standard benchmarks. If OpenEvolve-MARL can match or exceed
these results with open-source models --- and potentially improve upon
them via the fine-tuning loop --- the practical impact for practitioners
is immediate. CFR and PSRO are the backbone algorithms behind every
competitive poker AI, every deployed multi-agent negotiation system, and
much of the theoretical work on equilibrium computation. Faster, more
robust variants discovered by our system would be directly usable by the
community. Crucially, because our entire pipeline is open-source, any
novel algorithm we discover can be immediately inspected, verified, and
extended by others. Compare this to AlphaEvolve\textquotesingle s
discoveries, which are published as results but not as reproducible
artifacts --- our work would lower the barrier from
\textquotesingle interesting result\textquotesingle{} to
\textquotesingle usable contribution\textquotesingle{} to near zero.

\textbf{5.4 Positioning at the Frontier of Inference-Time Compute
Scaling}

The field is converging on a view that the most important axis of LLM
capability scaling in the near term is not parameters but inference-time
compute --- the ability to spend more FLOPs at test time to produce
better outputs. AdaEvolve is already an instance of this, using
evolutionary search as a structured inference-time procedure. Our
project pushes further: by fine-tuning the mutation model on search
outcomes, we are using inference-time experience to improve future
inference, closing a loop between test-time compute and model quality
that is novel in the algorithmic discovery literature. The 8 × H100
budget is well-matched to this: it is exactly the scale at which we can
run a high-throughput evolutionary search while simultaneously
fine-tuning a 7B model in the background, but it is far below the scale
at which proprietary systems operate. Demonstrating that the
self-improving loop yields measurable gains at this compute tier would
be a strong argument that the paradigm is accessible --- and that the
marginal return on open-source research infrastructure in this space is
high.

\textbf{5.5 Summary}

In short: if this project succeeds, it produces (a) an open,
reproducible implementation of the most advanced adaptive evolutionary
discovery system in the literature, (b) the first empirical evidence
that open-source LLMs with domain-specific fine-tuning can match
proprietary models in evolutionary algorithm search, (c) newly
discovered MARL algorithms that the community can immediately use and
build on, and (d) a general-purpose platform that other groups can
deploy for their own discovery tasks. Each of these outcomes is
independently valuable; together, they represent a contribution that is
disproportionately large relative to the compute investment required ---
which is precisely the kind of high-leverage research that 8 × H100 GPUs
are best suited to enable.

OpenEvolve-MARL \textbar{} Confidential Draft Proposal
